% Part: normal-modal-logic
% Chapter: axioms-systems
% Section: provability-properties

\documentclass[../../../include/open-logic-section]{subfiles}

\begin{document}

\olfileid[id]{nml}{prf}{prp}

\olsection{Sifat \usetoken{S}{derivability}}

\begin{prop}\ollabel{prop:derivabilityfacts}
  Misalkan $\Sigma$ suatu sistem modal dan $\Gamma$ suatu himpunan
  !!{formula}s modal. Sifat-sifat berikut berlaku:
  \begin{enumerate}
  \item\ollabel{prop:derivabilityfacts-monotonicity} \emph{Monotonisitas}:
    Jika $\Gamma\Proves[\Sigma]!A$ dan $\Gamma\subseteq\Delta$, maka
    $\Delta\Proves[\Sigma]!A$;
  \item\ollabel{prop:derivabilityfacts-reflexivity}
    \emph{Refleksivitas}: Jika $!A\in\Gamma$, maka
    $\Gamma\Proves[\Sigma]!A$;
  \item\ollabel{prop:derivabilityfacts-cut} \emph{Transitivitas}: Jika
    $\Gamma\Proves[\Sigma]!A$ dan
    $\Delta\cup\{!A\}\Proves[\Sigma]!B$, maka
    $\Gamma\cup\Delta\Proves[\Sigma]!B$;
  \item\ollabel{prop:derivabilityfacts-deduction} \emph{Teorema deduksi}:
    $\Gamma\cup\{!B\}\Proves[\Sigma]!A$ jika dan hanya jika
    $\Gamma\Proves[\Sigma]!B\lif !A$;
  \item\ollabel{prop:derivabilityfacts-ruleT} \emph{Aturan T}: Jika
    $\Gamma\Proves[\Sigma]!A_1$, \dots, $\Gamma\Proves[\Sigma]!A_n$, dan
    $!A_1\lif(!A_2\lif\cdots(!A_n\lif !B)\cdots)$ merupakan instansiasi
    tautologis, maka $\Gamma\Proves[\Sigma]!B$.
  \end{enumerate}
\end{prop}

Pembuktiannya merupakan latihan mudah. Bagian
\olref{prop:derivabilityfacts-ruleT} dari \olref{prop:derivabilityfacts}
memberi kita, misalnya, bahwa jika $\Gamma\Proves[\Sigma]!A\lor !B$ dan
$\Gamma\Proves[\Sigma]\lnot !A$, maka $\Gamma\Proves[\Sigma]!B$.
Selain itu, dalam pembahasan berikut kita menulis
$\Gamma,!A\Proves[\Sigma]!B$ alih-alih
$\Gamma\cup\{!A\}\Proves[\Sigma]!B$.

\begin{defn}
  Suatu himpunan~$\Gamma$ \emph{tertutup secara deduktif} relatif terhadap
  suatu sistem~$\Sigma$ jika dan hanya jika $\Gamma\Proves[\Sigma]!A$
  mengimplikasikan $!A\in\Gamma$.
\end{defn}

\end{document}
